\chapter{Reproducibility}
\label{app:reproducibility}

This appendix gives the information needed to inspect the reported analyses and distinguish a
repeatable code check from independent empirical replication. The full corpus and row-level
annotations are restricted. Detailed technical provenance is preserved in the project record;
this appendix focuses on the study design and evidence a reader needs to assess.

\section{Data provenance and analysis units}

The study uses the MIT Interview mock-interview distribution described by Naim et
al.~\cite{Naim2015}. Its later analysis established the interview-performance prediction
setting~\cite{Naim2018}. The present distribution contains 138 audiovisual sessions, 138
non-empty transcripts and one aggregate outcome per session. The 69 participants each have a
pre and post recording. A complete face and identity audit established the participant mapping
used by this thesis. It corrects the reciprocal post-session source assignments for participants
32 and 33 while leaving the raw files and their identifiers unchanged. Data joins use raw source
identifiers; folds and paired analyses use the corrected participant identifiers.

The primary outcome is the supplied aggregate \emph{RecommendHiring} rating. A raw rating on
the documented 1--7 scale is transformed to $[0,1]$ as $(y-1)/6$ before fitting. The raw
rating, aggregate row and transformed target are kept together for every session. Individual
crowd ratings are used only in the separate target-measurement audit. No job-performance or
employment outcome is available.

Before modelling, each of the 138 sessions passed the declared audio, video, transcript and
outcome joins. All 69 complete participant pairs met the three-modality eligibility rule. No
session was imputed or replaced with a legacy feature. The independent evaluation and resampling
unit is therefore the participant, not the session.

\section{Frozen representations}

Visual input is sampled from candidate-valid face crops selected by the one-hertz, centrally
located face rule in Chapter~\ref{ch:methodology}. At most 128 crops from each session enter
DINOv2-large. The complete face manifest covers every source session and all 17,664 selected
crops; the identity audit separately checked the participant-pair mapping. The positional face
rule is an operational candidate selector, not identity recognition.

Interview audio is the bounded 16-kHz mono signal from the same active interval as visual
sampling. It contains candidate and interviewer speech because the supplied channels do not
isolate speakers. WavLM-Base+ processes at most 16 selected non-overlapping 20-second windows.
Candidate-only transcript turns, selected by the audited role parser, are encoded with the
fixed all-mpnet-base-v2 sentence-transformer checkpoint. The corpus-wide text check found 77
retained and three excluded colonless fragments under the declared continuation rule.

All three encoders remain frozen. Their versions, preprocessing settings and output dimensions
were fixed for the accepted feature set. The same representations supply every fitted
configuration, so a comparison of
training objectives does not silently substitute another modality or checkpoint.

\section{Participant-safe fitting and model contrasts}

The fixed assignment places all 69 participant pairs into five folds of sizes 14, 14, 14, 14
and 13. Both sessions from one participant always share a fold. In each outer rotation, one
fold is held out for testing, the following fold is used for stopping, and the remaining
three are used for fitting. The assignment is deterministic and independent of outcomes and
perceived-attribute annotations; its exact recipe is retained in the provenance register.

Seven ordinary configurations were fitted at seeds 1, 11 and 21 in all five rotations:
ERM, visual-proxy sampling, interview-audio-proxy sampling, random-proxy W1,
visual-proxy W1, interview-audio-proxy W1 and ARL. These yield 105 learned fits.
Nine deterministic reference predictors use the same held-out roster but do not add seeded
fits. The separate perceived-gender-presentation oracle adds 15 fits after the ordinary
predictions are fixed. Every configuration--seed combination has one out-of-fold prediction
for each of the 138 sessions. Training-only scalers, proxy centres and stopping information
are not estimated from a held-out participant. The architecture, optimiser, sampler, penalty
and stopping rule are specified in Chapter~\ref{ch:methodology}.

The two principal W1 comparisons with ERM estimate the combined effect of proxy-balanced
sampling and the W1 penalty under this fixed design. Comparisons with the corresponding
sampling-only models isolate the incremental penalty. The random partition and ARL serve
different control questions. These contrasts do not assume that an unsupervised proxy group
is a demographic group.

\section{Perceived-attribute audit}

The accepted annotation source contains one observer's session-level judgements for all 69
post sessions and 57 pre sessions. It was supplied after ordinary model predictions and proxy
assignments were fixed. The ordinary model-fitting path had no access to those fields. The
exploratory oracle uses only the observed post-session F/M perceived-gender-presentation
codes for training participants in each rotation; held-out participant annotations remain
unavailable to the oracle. This field was selected after receipt of the source, so the
oracle is not a preregistered efficacy ceiling.

The primary audit joins the 69 post predictions to their corresponding post annotations.
The incomplete 57-session pre roster is a separate sensitivity, joined only to pre
predictions. A category requires at least ten distinct participants in its phase to
support a reported disparity estimate. Unknown and smaller categories are counted but not
merged to pass this rule. Perceived skin tone remains an ordinal 1--5 field, not a binary
attribute. The single-observer source lacks inter-rater agreement and cannot establish
participants' self-identified attributes.

\section{Statistical reproduction and interpretation}

For each original model and seed, correlations and errors are recomputed once from all five
concatenated held-out folds. Fold metrics are not averaged. A mean and sample standard
deviation across the three fixed seeds describe optimisation sensitivity, not uncertainty
over independent participants. The participant-cluster bootstrap samples whole participant
records with replacement 10,000 times at its fixed seed. For pooled utility, both sessions
remain together; for the primary audit, each sampled participant contributes the matching
post prediction and annotation. Intervals for the incomplete pre roster are computed
separately.

The maximum pairwise empirical W1 distance is calculated only among reportable categories
on each perceived-attribute axis. Within-seed configuration differences are formed first;
their arithmetic mean is then recomputed in each participant resample. Reported 95\%
percentile intervals are pointwise and are not simultaneous guarantees across all axes,
configurations and seeds. The target's own group distances and category-specific errors
are reported beside predicted-score disparity so that lower W1 is not mistaken for
predictive validity or hiring fairness.

\section{Exploratory linear follow-up}

The linear mechanism comparison reuses the accepted MIT features and participant folds,
not newly extracted representations. Its 15 deterministic configurations comprise a
natural task-only ridge reference, two proxy-balanced task-only references and visual
or audio W1 at strengths 0.002, 0.02 and 0.2 under both task exposures. Five outer
rotations yield 75 fits. Natural and balanced objectives use full-fold weights rather
than replaying the neural experiment's stochastic minibatches. The initial solution
is the exact task-only ridge optimum; the score output is linear and unclipped.

The empirical W1 integrates step quantiles exactly for unequal group sizes. Optimisation
uses the training design's row space and a fixed 1,000-iteration cap. Termination states
are retained: 62 fits converge and 13 highest-strength fits reach the cap. All 20 fits
underlying the four primary strength-0.02 comparisons converge. Pooled utility is
recomputed from concatenated OOF predictions, whereas proxy manipulation is the mean
of five fold-local W1 distances. Those summaries answer different questions and are
not interchangeable. Perceived-attribute contrasts use the same 69-post/57-pre source
and paired participant bootstrap as the original audit.

A separate outcome-only ridge-selection summary reports an unweighted mean of
five outer-fold metrics after participant-grouped inner selection. Its mean fold
correlation is not a pooled OOF correlation. The available summary does not supply
a paired participant interval against the fixed-alpha compact reference. This
descriptive result is therefore reported separately rather than used to declare
a statistically established improvement.

Proxy-association analysis also respects fold identities. Visual cluster IDs are aligned
using only shared training participants, before comparison with annotations. The
partitions agree up to permutation, permitting pooled own-test-fold association.
Audio partitions do not remain identical after relabelling and are reported fold-locally.
This analysis reuses frozen assignments and does not refit models or alter the
annotation-defined score disparities.

\section{Independent-corpus evidence}

AVI-Personality preserves its official participant splits: 452 training, 64 validation
and 130 test people. All 646 have six audio/video feature vectors; two missing training
outcomes reduce supervised fitting to 450. Equal-weight averaging across six question
clips precedes participant modelling. The visual encoder receives full frames rather
than MIT's face crops, and no text representation is supplied. These differences and
the original 1--5 recruiter-rating scale are part of the external analysis definition.

Featurewise scaling is fitted on training people only. Ridge alpha is selected by
validation MAE from the fixed five-value grid, without fitting on validation or test
participants. The final predictions are clipped to the source range. The independent
test baseline and self-reported-group MAEs are descriptive accepted summaries; no
paired utility-improvement interval or signed group-gap interval is supplied.
An absolute-gap percentile interval is not used as a signed disparity test.
The external baseline has no matched W1 intervention and must not be presented
as confirmation of mitigation.

\section{Seed-specific predictive utility}

Table~\ref{tab:stage-a-seeds} preserves every learned configuration's pooled held-out
utility at each fixed seed. These are repeated computations on the same 69 participants,
not 21 independent samples. Errors are in original 1--7 rating-scale units; normalised
errors equal these values divided by six, before display rounding.

\begin{table}[H]
\centering
\caption{Seed-specific pooled out-of-fold utility for the seven learned configurations. Each row
contains one metric calculation over all 138 held-out sessions; it is not an average of fold
metrics.}
\label{tab:stage-a-seeds}
\small
\renewcommand{\arraystretch}{1.05}
\begin{tabular}{llrrrr}
\toprule
Configuration & Seed & SRCC & Pearson & MAE (1--7) & RMSE (1--7) \\
\midrule
ERM & 1 & 0.110 & 0.110 & 0.926 & 1.151 \\
ERM & 11 & 0.089 & 0.093 & 1.189 & 1.488 \\
ERM & 21 & 0.311 & 0.296 & 0.737 & 0.922 \\
Face sampling & 1 & 0.078 & 0.071 & 0.977 & 1.218 \\
Face sampling & 11 & 0.174 & 0.191 & 0.990 & 1.223 \\
Face sampling & 21 & 0.325 & 0.310 & 0.779 & 0.954 \\
Audio sampling & 1 & 0.043 & 0.055 & 0.896 & 1.133 \\
Audio sampling & 11 & 0.108 & 0.133 & 0.965 & 1.191 \\
Audio sampling & 21 & 0.327 & 0.316 & 0.733 & 0.908 \\
Random W1 & 1 & 0.101 & 0.093 & 0.971 & 1.196 \\
Random W1 & 11 & 0.190 & 0.205 & 1.129 & 1.402 \\
Random W1 & 21 & 0.296 & 0.287 & 0.700 & 0.879 \\
Face W1 & 1 & 0.072 & 0.066 & 0.987 & 1.229 \\
Face W1 & 11 & 0.098 & 0.113 & 1.140 & 1.438 \\
Face W1 & 21 & 0.294 & 0.273 & 0.740 & 0.928 \\
Audio W1 & 1 & -0.007 & 0.006 & 0.969 & 1.199 \\
Audio W1 & 11 & 0.065 & 0.094 & 0.983 & 1.213 \\
Audio W1 & 21 & 0.280 & 0.253 & 0.728 & 0.909 \\
ARL & 1 & 0.072 & 0.082 & 0.908 & 1.130 \\
ARL & 11 & 0.087 & 0.093 & 1.193 & 1.490 \\
ARL & 21 & 0.312 & 0.298 & 0.738 & 0.923 \\
\bottomrule
\end{tabular}
\end{table}

\section{Evidence integrity and reader checks}

The accepted MIT production feature release completed all 138 sessions without a structured
failure. The accepted Stage-A release contains the 105 learned fits and nine deterministic
reference predictors; its exact resume check created no additional fit. Separate
annotation-free reports cover participant-bootstrap utility, proxy diagnostics, boundary
coverage, paired-change accuracy and inference-only modality masks. Stage B's 15 oracle
fits and the phase-aligned Stage-C aggregate were completed after the declared annotation
release. These outputs were frozen before this manuscript was revised. Earlier diagnostic
or preprocessing versions that failed the accepted boundary are not used as result
sources. The project record preserves the detailed version history needed to verify those
distinctions. The exploratory linear extension and independent AVI baseline remain
separate analyses with their own evidence and measurement boundaries; they do not
replace the original predictions.

A reader with only the manuscript can inspect the analysis units, fixed comparisons,
category counts, numerical tables and interpretation limits. A reader with the project
code can run its tests and rebuild the manuscript without raw media. Recomputing
participant-level empirical outputs requires authorised access to the original corpus,
annotation source, pinned encoders and frozen run inputs. The restricted participant-level
join and identifiable media are not publication artefacts; the aggregate Stage-C report
suffices to check the reported disparity tables and figure without rerunning the audit.
